% Part: first-order-logic
% Chapter: beyond
% Section: modal-logics

\documentclass[../../../include/open-logic-section]{subfiles}

\begin{document}

\olfileid{fol}{byd}{mod}

\olsection{मोडल तर्कशास्त्रे}

पुढील सशर्त वाक्याचे उदाहरण पाहा:
\begin{quote}
  जेरेमी त्या खोलीत एकटा असेल, तर त्याने मद्यपान केलेले आहे, तो विवस्त्र आहे
  आणि खुर्च्यांवर नाचत आहे.
\end{quote}
हे सशर्त विधान वास्तविक अभिव्यंजनाच्या अर्थाने सत्य असू शकते, पण तरीही
दिशाभूल करणारे असू शकते; कारण पूर्वांग असत्य आहे किंवा उत्तरांग सत्य आहे
एवढ्यापेक्षा पूर्वांग
आणि निष्कर्ष यांच्यात अधिक दृढ संबंध आहे असे ते सुचवताना दिसते. म्हणजे, हा
दावा फक्त या विशिष्ट जगातच सत्य नाही (जेरेमी खोलीत एकटा नसल्यामुळे येथे तो
क्षुल्लकपणे सत्य असू शकतो), तर पूर्वांग सत्य \emph{असते}, तर निष्कर्षही सत्य
\emph{असता}, असे त्याची शब्दरचना सुचवते. दुसऱ्या शब्दांत, हे विधान फक्त या
जगात नव्हे, तर कोणत्याही “संभाव्य” जगात सत्य आहे; किंवा ते या विशिष्ट जगात
केवळ सत्य असण्याऐवजी \emph{अनिवार्यपणे} सत्य आहे, असा त्याचा अर्थ घेता येतो.

अशा प्रकारच्या अनिवार्यतेचा अर्थ लावण्यासाठी मोडल तर्कशास्त्र रचले गेले.
नेहमीच्या विधानीय तर्कशास्त्रात चौकट-कारक जोडल्यावर मोडल विधानीय
तर्कशास्त्र मिळते; म्हणजे $!A$ हे !!a{formula} असेल, तर $\Box !A$ हेही
सूत्र असते. सहज अर्थाने, $\Box !A$ हे $!A$ \emph{अनिवार्यपणे} सत्य
आहे, किंवा कोणत्याही संभाव्य जगात सत्य आहे, असे सांगते. $\Diamond !A$ हे
सहसा $\lnot \Box \lnot !A$ चे संक्षिप्त रूप मानले जाते आणि $!A$
\emph{संभाव्यतः} सत्य आहे असे सांगणारे म्हणून वाचता येते. अर्थात, विधेय
तर्कशास्त्रातही मोडलता जोडता येते.

मोडल तर्कशास्त्राची चिन्हार्थमीमांसा देण्यासाठी क्रिप्के !!{structure}s वापरता
येतात; खरे तर क्रिप्के यांनी ही चिन्हार्थमीमांसा प्रथम मोडल तर्कशास्त्र
डोळ्यांसमोर ठेवूनच रचली. अंशतः क्रमांपुरते मर्यादित न राहता, अधिक सामान्यपणे
“संभाव्य जगांचा” $P$ हा संच आणि जगांमधील $\Atom{R}{x,y}$ हा द्विपदी
“प्राप्यता” संबंध घेतला जातो. सहज अर्थाने, $\Atom{R}{p,q}$ हे जग~$q$ हे
जग~$p$ शी सुसंगत आहे असे सांगते; म्हणजे आपण जग~$p$ “मध्ये” असू, तर जग
$q$ सारखे असू शकले असते ही शक्यता विचारात घ्यावी लागते.

मोडल तर्कशास्त्राला कधी कधी “विस्तारात्मक” तर्कशास्त्राच्या विरुद्ध
“आशयात्मक” तर्कशास्त्र म्हणतात. अभिजात तर्कशास्त्रासारख्या विस्तारात्मक
तर्कशास्त्राची अभिप्रेत चिन्हार्थमीमांसा केवळ “प्रत्यक्ष” अशा एका जगाचा
निर्देश करते; पण “आशयात्मक” तर्कशास्त्राची चिन्हार्थमीमांसा अधिक विस्तृत
सत्ताशास्त्रावर अवलंबून असते. अनिवार्यतेची !!{structure} रचण्याखेरीज,
उपयोजनानुसार $\Box$ आणि $\Diamond$ यांचा नवा अर्थ लावून इतर भाषिक रचनांची
!!{structure} रचण्यासाठीही मोडलता वापरता येते. उदाहरणार्थ:
\begin{enumerate}
\item सिद्धतायोग्यता तर्कशास्त्रात $\Box !A$ चा अर्थ “$!A$ सिद्धतायोग्य
  आहे” आणि $\Diamond !A$ चा अर्थ “$!A$ सुसंगत आहे” असा घेतला जातो.
\item ज्ञानविषयक तर्कशास्त्रात $\Box !A$ चा अर्थ “मला $!A$ माहीत आहे” किंवा
  “माझा $!A$ वर विश्वास आहे” असा घेता येतो.
\item कालिक तर्कशास्त्रात $\Box !A$ चा अर्थ “$!A$ नेहमी सत्य आहे” आणि
  $\Diamond !A$ चा अर्थ “$!A$ कधीतरी सत्य असते” असा घेता येतो.
\end{enumerate}

मोडलता हाताळणारे नियम आणि स्वयंसिद्धके तर्कशास्त्रात जोडावीशी वाटतात.
उदाहरणार्थ, $\Log{S4}$ प्रणालीत विधानीय तर्कशास्त्राची नेहमीची स्वयंसिद्धके
आणि नियम, तसेच पुढील स्वयंसिद्धके असतात:
\begin{align*}
& \Box (!A \lif !B) \lif (\Box !A \lif \Box !B)\\
& \Box !A \lif !A \\
& \Box !A \lif \Box \Box !A
\intertext{आणि “$!A$ पासून $\Box !A$ हा निष्कर्ष काढा” हा नियमही असतो.
  $\Log{S5}$ मध्ये पुढील स्वयंसिद्धक जोडले जाते:}
& \Diamond !A \lif \Box \Diamond !A
\end{align*}
या स्वयंसिद्धकांच्या विविध आवृत्त्या वेगवेगळ्या उपयोजनांसाठी योग्य असू शकतात;
उदाहरणार्थ, S5 हे सहसा तार्किक अनिवार्यतेची संकल्पना लक्षणित करते असे मानतात.
चांगली गोष्ट अशी की क्रिप्के !!{structure}s मधील प्राप्यता संबंध स्वाभाविक
रीतीने निर्बंधित करून, संबंधित !!{derivation} प्रणाली निर्दोष आणि संपूर्ण
असेल अशी चिन्हार्थमीमांसा सहसा सापडते. उदाहरणार्थ, $\Log{S4}$ अशा क्रिप्के
!!{structure}s च्या वर्गाशी संबंधित आहे ज्यांत प्राप्यता संबंध परावर्ती आणि
संक्रमक असतो. $\Log{S5}$ अशा क्रिप्के !!{structure}s च्या वर्गाशी संबंधित
आहे ज्यांत प्राप्यता संबंध {\em वैश्विक} असतो, म्हणजे प्रत्येक जग इतर प्रत्येक
जगापासून प्राप्य असते; म्हणून $\Box !A$ तेव्हा आणि केवळ तेव्हाच लागू होते,
जेव्हा $!A$ प्रत्येक जगात लागू होते.

\end{document}
